\chapter{Future Work}

Several directions remain open for future development of the pipeline.

\textbf{Cloud deployment.} The local Docker Compose stack is architecturally equivalent to a fully managed cloud deployment, but the migration itself has not been performed. Deploying the pipeline to Amazon MSK, Amazon Managed Service for Apache Flink, and Amazon RDS would validate that the architecture sustains the same latency and throughput guarantees under real network conditions and distributed resource allocation.

\textbf{Live sensor integration.} All evaluation in this work used replayed historical data. Integrating live sensor hardware — or a larger set of real-time API sources — would introduce network jitter, sensor dropout, and calibration drift that the replay mode cannot reproduce. Evaluating pipeline behaviour under these conditions, including the effect of the watermark on genuinely late-arriving readings, remains an important step toward production readiness.

\textbf{Fault tolerance.} The Flink job is configured with exactly-once semantics and periodic checkpointing, but recovery from node failure was not evaluated. Measuring cold-start recovery time — the delay between a TaskManager crash and the resumption of window processing from the last checkpoint — would quantify the practical availability of the pipeline under failure.

\textbf{Anomaly detection quality.} The current Z-score detector uses a static per-sensor baseline loaded at job startup. A rolling baseline updated from the live stream would adapt to gradual shifts in sensor output caused by seasonal variation or sensor degradation. More sophisticated approaches — including machine learning models trained on historical patterns — could reduce false positives further, particularly for stations near industrial sites with predictable periodic emission profiles.

\textbf{Broader geographic coverage.} The pipeline was designed and tested for 414 Ukrainian monitoring stations. Extending coverage to additional countries or integrating satellite-derived air quality estimates alongside ground-level sensors would increase the practical value of the system as a public health notification infrastructure.

\textbf{Targeted alert routing.} The current alerting consumer delivers all notifications to a single endpoint regardless of the geographic origin of the event. A more practical design would route alerts based on the location of the affected station — notifying regional health authorities, municipal services, or individual subscribers with a direct interest in a given area. This could be implemented through a publish-subscribe model in which consumers subscribe to alerts filtered by region or city, eliminating irrelevant notifications and reducing response latency for affected parties.